% Generated. Do not edit.
% Projection scriptural-date-location/section, version 1: the Scriptural Date and
% Location page, generated whole from research/chronology.toml, the Scripture
% chronology corpus and the source library by tools/tpt document-types refresh.
% document: liturgy/roman-rite/postconciliar/roman-missal-third-edition-en-us-2011/propers/temporal/pc-s54-twenty-eighth-sunday-in-ordinary-time-year-a
\runninghead{Scriptural Date and Location}
\section{Scriptural Date and Location}
\label{sec:date-location}

{\small Each passage the Mass reads or sings from Scripture, in canonical order. The Context column names what is dated --- the event the passage narrates, a setting it remembers, the traditional attribution of its author, or the composition of its book --- and where it belongs; the Date column gives the answer of the Scripture chronology corpus for that subject, with every alternative it records. Texts of the Mass that are not drawn from Scripture have no row. The evidence for each date follows the table.\par}
\medskip
% Generated. Do not edit.
% Deterministic projection of research/chronology.toml.
% schema: 5
% document: liturgy/roman-rite/postconciliar/roman-missal-third-edition-en-us-2011/propers/temporal/pc-s54-twenty-eighth-sunday-in-ordinary-time-year-a
% generated-by: tools/tpt proper-chronology annotations
\newcommand{\chronologyannotationclaim}[6]{#6}
\newcommand{\chronologyannotationreach}[3]{}
\newcommand{\chronologyannotationgroup}[3]{#3}
\newcommand{\chronologyannotation}[1]{%
  \ifcsname triptychchronologyannotation@#1\endcsname
    \csname triptychchronologyannotation@#1\endcsname
  \else
    \PackageError{triptych}{No chronology annotation for #1}{}%
  \fi}
% appointment-dependency: guidance/liturgy/postconciliar-propers-registry.md sha256=376860a7c37a4c75bb6d880cd1b260183be1877a07e3cd01660a9e2854e2ea72
% appointment-dependency: src/claude/liturgy/roman-rite/postconciliar/roman-missal-third-edition-en-us-2011/propers/registry/README.md sha256=22626b14e0506b69c0ddc254b0ade2b6f7c9ae7cf8ef98fe356746a8f47fb1eb
% appointment-dependency: src/claude/liturgy/roman-rite/postconciliar/roman-missal-third-edition-en-us-2011/propers/registry/formula-dispositions.md sha256=37735c7edc0fbf692a691c7c4b0487b80831108daf91621b91f7b8858e190879
% appointment-dependency: src/claude/liturgy/roman-rite/postconciliar/roman-missal-third-edition-en-us-2011/propers/temporal/pc-s54-twenty-eighth-sunday-in-ordinary-time-year-a/instance/manifest.md sha256=5f9a0a20204bb48e1f099cc2ead60e603dd13b27568f5e1ab8febf6bc7e3ccbb
% appointment-dependency: src/claude/liturgy/roman-rite/postconciliar/roman-missal-third-edition-en-us-2011/propers/temporal/pc-s54-twenty-eighth-sunday-in-ordinary-time-year-a/propers/verified.md sha256=5e5943cf75a72febf05fe6c61c8b1b028a15ba5f94d1b7e90d7d945e369b8984
% appointment-dependency: src/claude/liturgy/roman-rite/postconciliar/roman-missal-third-edition-en-us-2011/propers/temporal/pc-s54-twenty-eighth-sunday-in-ordinary-time-year-a/research/chronology-inputs.toml sha256=32b90d0e0f0e9513a951e89e4cde71e5c09df1375276fa1fef3cde76fd5c8181
% appointment-dependency: src/claude/liturgy/roman-rite/postconciliar/roman-missal-third-edition-en-us-2011/propers/temporal/shared/ordinary-time/weeks/28/propers/verified.md sha256=013744f84531e00098f06f8bd60f097502717ec3fef2e321af20beab074f6342
% appointment-note: entrance-antiphon: Psalm 129:3-4; adaptation; vulgate numbering; canonical verse queries, not appointed wording.
% appointment-note: collect: Composed oration with no scriptural locator; its vocabulary of grace going before and following is a loose echo of no one verse, which the page-2 rule excludes.
% appointment-note: first-reading: Isaiah 25:6-10a; appointed; vulgate numbering; whole-verse query envelope, not subverse chronology or appointed wording.
% appointment-note: responsorial-psalm: Psalm 23:1-3a, 3b-4, 5, 6; appointed; hebrew numbering; whole-verse query envelope, not subverse chronology or appointed wording.
% appointment-note: responsorial-psalm: Psalm 23:6cd; response; hebrew numbering; whole-verse query envelope, not subverse chronology or appointed wording.
% appointment-note: second-reading: Philippians 4:12-14, 19-20; appointed; vulgate numbering; canonical verse queries, not appointed wording.
% appointment-note: gospel-acclamation: Cf. Ephesians 1:17-18; adaptation; vulgate numbering; canonical verse queries, not appointed wording.
% appointment-note: gospel-long: Matthew 22:1-14; appointed; vulgate numbering; canonical verse queries, not appointed wording.
% appointment-note: gospel-short: Matthew 22:1-10; appointed; vulgate numbering; canonical verse queries, not appointed wording.
% appointment-note: prayer-over-offerings: Composed oration with no scriptural locator.
% appointment-note: communion-antiphon-a: Cf. Psalm 33:11; adaptation; vulgate numbering; canonical verse queries, not appointed wording.
% appointment-note: communion-antiphon-b: 1 John 3:2; adaptation; vulgate numbering; canonical verse queries, not appointed wording.
% appointment-note: prayer-after-communion: Composed oration with no scriptural locator. Its words divinae naturae consortes also stand in 2 Peter 1:4 in the Clementine; that is a loose echo the Missal does not cite, which the page-2 rule excludes.
% entrance-antiphon: Entrance Antiphon
\expandafter\def\csname triptychchronologyannotation@entrance-antiphon\endcsname{%
  \chronologyannotationgroup{historical-setting}{research-pending}{\textbf{Historical setting} -- Occasion date unresolved.}%
  \space \chronologyannotationgroup{composition}{preferred}{\textbf{Composition}: \chronologyannotationreach{Ps.129.3}{Ps}{true}\chronologyannotationreach{Ps.129.4}{Ps}{true}\chronologyannotationclaim{critical.psalms.latest-composition-boundary}{composition}{catholic-critical-v1}{preferred}{before the Maccabean period, around 165 B.C.; no individual psalm can be dated securely}{Before c. 165 B.C.}}%
}
% first-reading: First Reading
\expandafter\def\csname triptychchronologyannotation@first-reading\endcsname{%
  \chronologyannotationgroup{traditional-attribution}{preferred}{\textbf{Traditional attribution}: \chronologyannotationreach{Is.25.10}{Is}{true}\chronologyannotationreach{Is.25.6}{Is}{true}\chronologyannotationreach{Is.25.7}{Is}{true}\chronologyannotationreach{Is.25.8}{Is}{true}\chronologyannotationreach{Is.25.9}{Is}{true}\chronologyannotationclaim{israel.prophets.isaias-traditional-era}{traditional-attribution}{catholic-traditional-v1}{preferred}{from the closing year of Ozias, King of Juda (740 B.C.); no prophecies appear to be later than 701 B.C.}{Isaias (ministry in Souvay's traditional account), B.C. 740--701}.}%
}
% responsorial-psalm: Responsorial Psalm
\expandafter\def\csname triptychchronologyannotation@responsorial-psalm\endcsname{%
  \chronologyannotationgroup{historical-setting}{research-pending}{\textbf{Historical setting} -- Occasion date unresolved.}%
  \space \chronologyannotationgroup{composition}{preferred}{\textbf{Composition}: \chronologyannotationreach{Ps.22.1}{Ps}{true}\chronologyannotationreach{Ps.22.2}{Ps}{true}\chronologyannotationreach{Ps.22.3}{Ps}{true}\chronologyannotationreach{Ps.22.4}{Ps}{true}\chronologyannotationreach{Ps.22.5}{Ps}{true}\chronologyannotationreach{Ps.22.6}{Ps}{true}\chronologyannotationclaim{critical.psalms.latest-composition-boundary}{composition}{catholic-critical-v1}{preferred}{before the Maccabean period, around 165 B.C.; no individual psalm can be dated securely}{Before c. 165 B.C.}}%
}
% second-reading: Second Reading
\expandafter\def\csname triptychchronologyannotation@second-reading\endcsname{%
  \chronologyannotationgroup{composition}{disputed}{\textbf{Composition} -- disputed: \chronologyannotationreach{Phil.4.12}{Phil}{true}\chronologyannotationreach{Phil.4.13}{Phil}{true}\chronologyannotationreach{Phil.4.14}{Phil}{true}\chronologyannotationreach{Phil.4.19}{Phil}{true}\chronologyannotationreach{Phil.4.20}{Phil}{true}\chronologyannotationclaim{composition.epistle-to-the-philippians}{composition}{catholic-traditional-v1}{disputed}{(Philemon; Colossians; Ephesians; Philippians), 61}{A.D. 61}; \chronologyannotationreach{Phil.4.12}{Phil}{true}\chronologyannotationreach{Phil.4.13}{Phil}{true}\chronologyannotationreach{Phil.4.14}{Phil}{true}\chronologyannotationreach{Phil.4.19}{Phil}{true}\chronologyannotationreach{Phil.4.20}{Phil}{true}\chronologyannotationclaim{composition.epistle-to-the-philippians}{composition}{catholic-traditional-v1}{disputed}{at Rome (A.D. 62-64)}{A.D. 62--64}.}%
}
% gospel-acclamation: Alleluia Verse
\expandafter\def\csname triptychchronologyannotation@gospel-acclamation\endcsname{%
  \chronologyannotationgroup{composition}{disputed}{\textbf{Composition} -- disputed: \chronologyannotationreach{Eph.1.17}{Eph}{true}\chronologyannotationreach{Eph.1.18}{Eph}{true}\chronologyannotationclaim{composition.epistle-to-the-ephesians}{composition}{catholic-traditional-v1}{disputed}{a period between 58 and 63}{A.D. 58--63}; \chronologyannotationreach{Eph.1.17}{Eph}{true}\chronologyannotationreach{Eph.1.18}{Eph}{true}\chronologyannotationclaim{composition.epistle-to-the-ephesians}{composition}{catholic-traditional-v1}{disputed}{(Philemon; Colossians; Ephesians; Philippians), 61}{A.D. 61}.}%
}
% gospel-long: Gospel (longer form)
\expandafter\def\csname triptychchronologyannotation@gospel-long\endcsname{%
  \chronologyannotationgroup{narrated-event}{preferred}{\textbf{Event}: \chronologyannotationreach{Matt.22.1}{Matt.22.1-14}{false}\chronologyannotationreach{Matt.22.10}{Matt.22.1-14}{false}\chronologyannotationreach{Matt.22.11}{Matt.22.1-14}{false}\chronologyannotationreach{Matt.22.12}{Matt.22.1-14}{false}\chronologyannotationreach{Matt.22.13}{Matt.22.1-14}{false}\chronologyannotationreach{Matt.22.14}{Matt.22.1-14}{false}\chronologyannotationreach{Matt.22.2}{Matt.22.1-14}{false}\chronologyannotationreach{Matt.22.3}{Matt.22.1-14}{false}\chronologyannotationreach{Matt.22.4}{Matt.22.1-14}{false}\chronologyannotationreach{Matt.22.5}{Matt.22.1-14}{false}\chronologyannotationreach{Matt.22.6}{Matt.22.1-14}{false}\chronologyannotationreach{Matt.22.7}{Matt.22.1-14}{false}\chronologyannotationreach{Matt.22.8}{Matt.22.1-14}{false}\chronologyannotationreach{Matt.22.9}{Matt.22.1-14}{false}\chronologyannotationclaim{life-of-christ.parable-of-the-marriage-feast}{narrated-event}{catholic-traditional-v1}{preferred}{derived A.D. 29 from February, A.U.C. 782- Passover, 782}{The parable of the marriage feast, proposed in the Temple, A.D. 29 (derived)}.}%
  \space \chronologyannotationgroup{composition}{disputed}{\textbf{Composition} -- disputed: \chronologyannotationreach{Matt.22.1}{Matt}{true}\chronologyannotationreach{Matt.22.10}{Matt}{true}\chronologyannotationreach{Matt.22.11}{Matt}{true}\chronologyannotationreach{Matt.22.12}{Matt}{true}\chronologyannotationreach{Matt.22.13}{Matt}{true}\chronologyannotationreach{Matt.22.14}{Matt}{true}\chronologyannotationreach{Matt.22.2}{Matt}{true}\chronologyannotationreach{Matt.22.3}{Matt}{true}\chronologyannotationreach{Matt.22.4}{Matt}{true}\chronologyannotationreach{Matt.22.5}{Matt}{true}\chronologyannotationreach{Matt.22.6}{Matt}{true}\chronologyannotationreach{Matt.22.7}{Matt}{true}\chronologyannotationreach{Matt.22.8}{Matt}{true}\chronologyannotationreach{Matt.22.9}{Matt}{true}\chronologyannotationclaim{composition.gospel-of-matthew}{composition}{catholic-traditional-v1}{disputed}{about A.D. 38-45}{c. A.D. 38--45}; \chronologyannotationreach{Matt.22.1}{Matt}{true}\chronologyannotationreach{Matt.22.10}{Matt}{true}\chronologyannotationreach{Matt.22.11}{Matt}{true}\chronologyannotationreach{Matt.22.12}{Matt}{true}\chronologyannotationreach{Matt.22.13}{Matt}{true}\chronologyannotationreach{Matt.22.14}{Matt}{true}\chronologyannotationreach{Matt.22.2}{Matt}{true}\chronologyannotationreach{Matt.22.3}{Matt}{true}\chronologyannotationreach{Matt.22.4}{Matt}{true}\chronologyannotationreach{Matt.22.5}{Matt}{true}\chronologyannotationreach{Matt.22.6}{Matt}{true}\chronologyannotationreach{Matt.22.7}{Matt}{true}\chronologyannotationreach{Matt.22.8}{Matt}{true}\chronologyannotationreach{Matt.22.9}{Matt}{true}\chronologyannotationclaim{composition.gospel-of-matthew}{composition}{catholic-traditional-v1}{disputed}{about the year 40-42}{c. A.D. 40--42}; \chronologyannotationreach{Matt.22.1}{Matt}{true}\chronologyannotationreach{Matt.22.10}{Matt}{true}\chronologyannotationreach{Matt.22.11}{Matt}{true}\chronologyannotationreach{Matt.22.12}{Matt}{true}\chronologyannotationreach{Matt.22.13}{Matt}{true}\chronologyannotationreach{Matt.22.14}{Matt}{true}\chronologyannotationreach{Matt.22.2}{Matt}{true}\chronologyannotationreach{Matt.22.3}{Matt}{true}\chronologyannotationreach{Matt.22.4}{Matt}{true}\chronologyannotationreach{Matt.22.5}{Matt}{true}\chronologyannotationreach{Matt.22.6}{Matt}{true}\chronologyannotationreach{Matt.22.7}{Matt}{true}\chronologyannotationreach{Matt.22.8}{Matt}{true}\chronologyannotationreach{Matt.22.9}{Matt}{true}\chronologyannotationclaim{composition.gospel-of-matthew}{composition}{catholic-traditional-v1}{disputed}{the years 40-45}{A.D. 40--45}; \chronologyannotationreach{Matt.22.1}{Matt}{true}\chronologyannotationreach{Matt.22.10}{Matt}{true}\chronologyannotationreach{Matt.22.11}{Matt}{true}\chronologyannotationreach{Matt.22.12}{Matt}{true}\chronologyannotationreach{Matt.22.13}{Matt}{true}\chronologyannotationreach{Matt.22.14}{Matt}{true}\chronologyannotationreach{Matt.22.2}{Matt}{true}\chronologyannotationreach{Matt.22.3}{Matt}{true}\chronologyannotationreach{Matt.22.4}{Matt}{true}\chronologyannotationreach{Matt.22.5}{Matt}{true}\chronologyannotationreach{Matt.22.6}{Matt}{true}\chronologyannotationreach{Matt.22.7}{Matt}{true}\chronologyannotationreach{Matt.22.8}{Matt}{true}\chronologyannotationreach{Matt.22.9}{Matt}{true}\chronologyannotationclaim{composition.gospel-of-matthew}{composition}{catholic-traditional-v1}{disputed}{about the year 60-68}{c. A.D. 60--68}; \chronologyannotationreach{Matt.22.1}{Matt}{true}\chronologyannotationreach{Matt.22.10}{Matt}{true}\chronologyannotationreach{Matt.22.11}{Matt}{true}\chronologyannotationreach{Matt.22.12}{Matt}{true}\chronologyannotationreach{Matt.22.13}{Matt}{true}\chronologyannotationreach{Matt.22.14}{Matt}{true}\chronologyannotationreach{Matt.22.2}{Matt}{true}\chronologyannotationreach{Matt.22.3}{Matt}{true}\chronologyannotationreach{Matt.22.4}{Matt}{true}\chronologyannotationreach{Matt.22.5}{Matt}{true}\chronologyannotationreach{Matt.22.6}{Matt}{true}\chronologyannotationreach{Matt.22.7}{Matt}{true}\chronologyannotationreach{Matt.22.8}{Matt}{true}\chronologyannotationreach{Matt.22.9}{Matt}{true}\chronologyannotationclaim{composition.gospel-of-matthew}{composition}{catholic-traditional-v1}{disputed}{about the years 64-67}{c. A.D. 64--67}; \chronologyannotationreach{Matt.22.1}{Matt}{true}\chronologyannotationreach{Matt.22.10}{Matt}{true}\chronologyannotationreach{Matt.22.11}{Matt}{true}\chronologyannotationreach{Matt.22.12}{Matt}{true}\chronologyannotationreach{Matt.22.13}{Matt}{true}\chronologyannotationreach{Matt.22.14}{Matt}{true}\chronologyannotationreach{Matt.22.2}{Matt}{true}\chronologyannotationreach{Matt.22.3}{Matt}{true}\chronologyannotationreach{Matt.22.4}{Matt}{true}\chronologyannotationreach{Matt.22.5}{Matt}{true}\chronologyannotationreach{Matt.22.6}{Matt}{true}\chronologyannotationreach{Matt.22.7}{Matt}{true}\chronologyannotationreach{Matt.22.8}{Matt}{true}\chronologyannotationreach{Matt.22.9}{Matt}{true}\chronologyannotationclaim{composition.gospel-of-matthew}{composition}{catholic-traditional-v1}{disputed}{about the year 50}{c. A.D. 50}.}%
}
% gospel-short: Gospel (shorter form)
\expandafter\def\csname triptychchronologyannotation@gospel-short\endcsname{%
  \chronologyannotationgroup{narrated-event}{preferred}{\textbf{Event}: \chronologyannotationreach{Matt.22.1}{Matt.22.1-14}{false}\chronologyannotationreach{Matt.22.10}{Matt.22.1-14}{false}\chronologyannotationreach{Matt.22.2}{Matt.22.1-14}{false}\chronologyannotationreach{Matt.22.3}{Matt.22.1-14}{false}\chronologyannotationreach{Matt.22.4}{Matt.22.1-14}{false}\chronologyannotationreach{Matt.22.5}{Matt.22.1-14}{false}\chronologyannotationreach{Matt.22.6}{Matt.22.1-14}{false}\chronologyannotationreach{Matt.22.7}{Matt.22.1-14}{false}\chronologyannotationreach{Matt.22.8}{Matt.22.1-14}{false}\chronologyannotationreach{Matt.22.9}{Matt.22.1-14}{false}\chronologyannotationclaim{life-of-christ.parable-of-the-marriage-feast}{narrated-event}{catholic-traditional-v1}{preferred}{derived A.D. 29 from February, A.U.C. 782- Passover, 782}{The parable of the marriage feast, proposed in the Temple, A.D. 29 (derived)}.}%
  \space \chronologyannotationgroup{composition}{disputed}{\textbf{Composition} -- disputed: \chronologyannotationreach{Matt.22.1}{Matt}{true}\chronologyannotationreach{Matt.22.10}{Matt}{true}\chronologyannotationreach{Matt.22.2}{Matt}{true}\chronologyannotationreach{Matt.22.3}{Matt}{true}\chronologyannotationreach{Matt.22.4}{Matt}{true}\chronologyannotationreach{Matt.22.5}{Matt}{true}\chronologyannotationreach{Matt.22.6}{Matt}{true}\chronologyannotationreach{Matt.22.7}{Matt}{true}\chronologyannotationreach{Matt.22.8}{Matt}{true}\chronologyannotationreach{Matt.22.9}{Matt}{true}\chronologyannotationclaim{composition.gospel-of-matthew}{composition}{catholic-traditional-v1}{disputed}{about A.D. 38-45}{c. A.D. 38--45}; \chronologyannotationreach{Matt.22.1}{Matt}{true}\chronologyannotationreach{Matt.22.10}{Matt}{true}\chronologyannotationreach{Matt.22.2}{Matt}{true}\chronologyannotationreach{Matt.22.3}{Matt}{true}\chronologyannotationreach{Matt.22.4}{Matt}{true}\chronologyannotationreach{Matt.22.5}{Matt}{true}\chronologyannotationreach{Matt.22.6}{Matt}{true}\chronologyannotationreach{Matt.22.7}{Matt}{true}\chronologyannotationreach{Matt.22.8}{Matt}{true}\chronologyannotationreach{Matt.22.9}{Matt}{true}\chronologyannotationclaim{composition.gospel-of-matthew}{composition}{catholic-traditional-v1}{disputed}{about the year 40-42}{c. A.D. 40--42}; \chronologyannotationreach{Matt.22.1}{Matt}{true}\chronologyannotationreach{Matt.22.10}{Matt}{true}\chronologyannotationreach{Matt.22.2}{Matt}{true}\chronologyannotationreach{Matt.22.3}{Matt}{true}\chronologyannotationreach{Matt.22.4}{Matt}{true}\chronologyannotationreach{Matt.22.5}{Matt}{true}\chronologyannotationreach{Matt.22.6}{Matt}{true}\chronologyannotationreach{Matt.22.7}{Matt}{true}\chronologyannotationreach{Matt.22.8}{Matt}{true}\chronologyannotationreach{Matt.22.9}{Matt}{true}\chronologyannotationclaim{composition.gospel-of-matthew}{composition}{catholic-traditional-v1}{disputed}{the years 40-45}{A.D. 40--45}; \chronologyannotationreach{Matt.22.1}{Matt}{true}\chronologyannotationreach{Matt.22.10}{Matt}{true}\chronologyannotationreach{Matt.22.2}{Matt}{true}\chronologyannotationreach{Matt.22.3}{Matt}{true}\chronologyannotationreach{Matt.22.4}{Matt}{true}\chronologyannotationreach{Matt.22.5}{Matt}{true}\chronologyannotationreach{Matt.22.6}{Matt}{true}\chronologyannotationreach{Matt.22.7}{Matt}{true}\chronologyannotationreach{Matt.22.8}{Matt}{true}\chronologyannotationreach{Matt.22.9}{Matt}{true}\chronologyannotationclaim{composition.gospel-of-matthew}{composition}{catholic-traditional-v1}{disputed}{about the year 60-68}{c. A.D. 60--68}; \chronologyannotationreach{Matt.22.1}{Matt}{true}\chronologyannotationreach{Matt.22.10}{Matt}{true}\chronologyannotationreach{Matt.22.2}{Matt}{true}\chronologyannotationreach{Matt.22.3}{Matt}{true}\chronologyannotationreach{Matt.22.4}{Matt}{true}\chronologyannotationreach{Matt.22.5}{Matt}{true}\chronologyannotationreach{Matt.22.6}{Matt}{true}\chronologyannotationreach{Matt.22.7}{Matt}{true}\chronologyannotationreach{Matt.22.8}{Matt}{true}\chronologyannotationreach{Matt.22.9}{Matt}{true}\chronologyannotationclaim{composition.gospel-of-matthew}{composition}{catholic-traditional-v1}{disputed}{about the years 64-67}{c. A.D. 64--67}; \chronologyannotationreach{Matt.22.1}{Matt}{true}\chronologyannotationreach{Matt.22.10}{Matt}{true}\chronologyannotationreach{Matt.22.2}{Matt}{true}\chronologyannotationreach{Matt.22.3}{Matt}{true}\chronologyannotationreach{Matt.22.4}{Matt}{true}\chronologyannotationreach{Matt.22.5}{Matt}{true}\chronologyannotationreach{Matt.22.6}{Matt}{true}\chronologyannotationreach{Matt.22.7}{Matt}{true}\chronologyannotationreach{Matt.22.8}{Matt}{true}\chronologyannotationreach{Matt.22.9}{Matt}{true}\chronologyannotationclaim{composition.gospel-of-matthew}{composition}{catholic-traditional-v1}{disputed}{about the year 50}{c. A.D. 50}.}%
}
% communion-antiphon-a: Communion Antiphon (first)
\expandafter\def\csname triptychchronologyannotation@communion-antiphon-a\endcsname{%
  \chronologyannotationgroup{traditional-attribution}{disputed}{\textbf{Traditional attribution} -- disputed: \chronologyannotationreach{Ps.33.11}{Ps.33 Ps.39 Ps.70 Ps.85 Ps.94 Ps.97}{true}\chronologyannotationclaim{israel.monarchy.david-traditional-era}{traditional-attribution}{catholic-traditional-v1}{disputed}{reigned from 1055 to 1015 B.C.}{David (reign in the usual chronology), B.C. 1055--1015}.}%
  \space \chronologyannotationgroup{superscription-setting}{preferred}{\textbf{Superscription setting}: \chronologyannotationreach{Ps.33.11}{Ps.33 Ps.55}{true}\chronologyannotationclaim{israel.monarchy.david-at-geth}{superscription-setting}{catholic-traditional-v1}{preferred}{A. M. 2944, A. C. 1060}{David's flight to Achis, king of Geth, A.M. 2944}.}%
  \space \chronologyannotationgroup{composition}{preferred}{\textbf{Composition}: \chronologyannotationreach{Ps.33.11}{Ps}{true}\chronologyannotationclaim{critical.psalms.latest-composition-boundary}{composition}{catholic-critical-v1}{preferred}{before the Maccabean period, around 165 B.C.; no individual psalm can be dated securely}{Before c. 165 B.C.}}%
}
% communion-antiphon-b: Communion Antiphon (alternative)
\expandafter\def\csname triptychchronologyannotation@communion-antiphon-b\endcsname{%
  \chronologyannotationgroup{composition}{preferred}{\textbf{Composition}: \chronologyannotationreach{1John.3.2}{1John}{true}\chronologyannotationclaim{composition.first-epistle-of-st-john}{composition}{catholic-traditional-v1}{preferred}{from the year 90 to 100 (approximately)}{c. A.D. 90--100}.}%
}
% presentation: orientation-v3
\newcommand{\chronologyorientationrow}[2]{%
  \ifcsname triptychchronologyorientation@#1\endcsname
    \csname triptychchronologyorientation@#1\endcsname{#2}%
  \else
    \PackageError{triptych}{No chronology orientation row for #1}{}%
  \fi}
\expandafter\def\csname triptychchronologyorientation@row-ee4de941fd327b98\endcsname#1{%
Responsorial Psalm: Ps. 23:1--3a, 3b--4, 5, 6 (Heb.; Vg. 22), Ps. 23:6cd (Heb.; Vg. 22) & Historical setting & Occasion date unresolved\\*[.1em]%
}
\expandafter\def\csname triptychchronologyorientation@row-163cb89db43056f7\endcsname#1{%
 & Composition: Psalms & \chronologyannotationreach{Ps.22.1}{Ps}{true}\chronologyannotationreach{Ps.22.2}{Ps}{true}\chronologyannotationreach{Ps.22.3}{Ps}{true}\chronologyannotationreach{Ps.22.4}{Ps}{true}\chronologyannotationreach{Ps.22.5}{Ps}{true}\chronologyannotationreach{Ps.22.6}{Ps}{true}\chronologyannotationclaim{critical.psalms.latest-composition-boundary}{composition}{catholic-critical-v1}{preferred}{before the Maccabean period, around 165 B.C.; no individual psalm can be dated securely}{Before c.~165 B.C.}\\[.3em]%
}
\expandafter\def\csname triptychchronologyorientation@row-282b657f2d8824fb\endcsname#1{%
Communion Antiphon (first): Cf. Psalm 33:11 (Vg.; Heb. 34) & Traditional attribution: David's reign; Hebron, then Jerusalem & Disputed: \chronologyannotationreach{Ps.33.11}{Ps.33 Ps.39 Ps.70 Ps.85 Ps.94 Ps.97}{true}\chronologyannotationclaim{israel.monarchy.david-traditional-era}{traditional-attribution}{catholic-traditional-v1}{disputed}{reigned from 1055 to 1015 B.C.}{B.C.~1055--1015}\\*[.1em]%
}
\expandafter\def\csname triptychchronologyorientation@row-2b64174a6ff30e1b\endcsname#1{%
 & Superscription setting: David's flight to Achis, king of Geth; Geth, at the Philistine court & \chronologyannotationreach{Ps.33.11}{Ps.33 Ps.55}{true}\chronologyannotationclaim{israel.monarchy.david-at-geth}{superscription-setting}{catholic-traditional-v1}{preferred}{A. M. 2944, A. C. 1060}{A.M.~2944 (Haydock edition reporting Ussher)}\\*[.1em]%
}
\expandafter\def\csname triptychchronologyorientation@row-97d3f578b255df94\endcsname#1{%
 & Composition: Psalms & \chronologyannotationreach{Ps.33.11}{Ps}{true}\chronologyannotationclaim{critical.psalms.latest-composition-boundary}{composition}{catholic-critical-v1}{preferred}{before the Maccabean period, around 165 B.C.; no individual psalm can be dated securely}{Before c.~165 B.C.}\\[.3em]%
}
\expandafter\def\csname triptychchronologyorientation@row-0c09cf47b3a0c749\endcsname#1{%
Entrance Antiphon: Ps. 129:3--4 (Vg.; Heb. 130) & Historical setting & Occasion date unresolved\\*[.1em]%
}
\expandafter\def\csname triptychchronologyorientation@row-410350b07d43c785\endcsname#1{%
 & Composition: Psalms & \chronologyannotationreach{Ps.129.3}{Ps}{true}\chronologyannotationreach{Ps.129.4}{Ps}{true}\chronologyannotationclaim{critical.psalms.latest-composition-boundary}{composition}{catholic-critical-v1}{preferred}{before the Maccabean period, around 165 B.C.; no individual psalm can be dated securely}{Before c.~165 B.C.}\\[.3em]%
}
\expandafter\def\csname triptychchronologyorientation@row-b81bab3a554aee2e\endcsname#1{%
First Reading: Is. 25:6--10a & Traditional attribution: Isaias (ministry in Souvay's traditional account); Judah and Jerusalem & \chronologyannotationreach{Is.25.10}{Is}{true}\chronologyannotationreach{Is.25.6}{Is}{true}\chronologyannotationreach{Is.25.7}{Is}{true}\chronologyannotationreach{Is.25.8}{Is}{true}\chronologyannotationreach{Is.25.9}{Is}{true}\chronologyannotationclaim{israel.prophets.isaias-traditional-era}{traditional-attribution}{catholic-traditional-v1}{preferred}{from the closing year of Ozias, King of Juda (740 B.C.); no prophecies appear to be later than 701 B.C.}{B.C.~740--701}\\[.3em]%
}
\expandafter\def\csname triptychchronologyorientation@row-bbdb0e76aa03a178\endcsname#1{%
Gospel (longer form): Matt. 22:1--14 & Narrated event: The parable of the marriage feast, proposed in the Temple & \chronologyannotationreach{Matt.22.1}{Matt.22.1-14}{false}\chronologyannotationreach{Matt.22.10}{Matt.22.1-14}{false}\chronologyannotationreach{Matt.22.11}{Matt.22.1-14}{false}\chronologyannotationreach{Matt.22.12}{Matt.22.1-14}{false}\chronologyannotationreach{Matt.22.13}{Matt.22.1-14}{false}\chronologyannotationreach{Matt.22.14}{Matt.22.1-14}{false}\chronologyannotationreach{Matt.22.2}{Matt.22.1-14}{false}\chronologyannotationreach{Matt.22.3}{Matt.22.1-14}{false}\chronologyannotationreach{Matt.22.4}{Matt.22.1-14}{false}\chronologyannotationreach{Matt.22.5}{Matt.22.1-14}{false}\chronologyannotationreach{Matt.22.6}{Matt.22.1-14}{false}\chronologyannotationreach{Matt.22.7}{Matt.22.1-14}{false}\chronologyannotationreach{Matt.22.8}{Matt.22.1-14}{false}\chronologyannotationreach{Matt.22.9}{Matt.22.1-14}{false}\chronologyannotationclaim{life-of-christ.parable-of-the-marriage-feast}{narrated-event}{catholic-traditional-v1}{preferred}{derived A.D. 29 from February, A.U.C. 782- Passover, 782}{A.D.~29 (derived)}\\*[.1em]%
}
\expandafter\def\csname triptychchronologyorientation@row-4d9bb5eddc53b491\endcsname#1{%
 & Composition: Matthew; Semitic original associated with Palestine & Disputed: \chronologyannotationreach{Matt.22.1}{Matt}{true}\chronologyannotationreach{Matt.22.10}{Matt}{true}\chronologyannotationreach{Matt.22.11}{Matt}{true}\chronologyannotationreach{Matt.22.12}{Matt}{true}\chronologyannotationreach{Matt.22.13}{Matt}{true}\chronologyannotationreach{Matt.22.14}{Matt}{true}\chronologyannotationreach{Matt.22.2}{Matt}{true}\chronologyannotationreach{Matt.22.3}{Matt}{true}\chronologyannotationreach{Matt.22.4}{Matt}{true}\chronologyannotationreach{Matt.22.5}{Matt}{true}\chronologyannotationreach{Matt.22.6}{Matt}{true}\chronologyannotationreach{Matt.22.7}{Matt}{true}\chronologyannotationreach{Matt.22.8}{Matt}{true}\chronologyannotationreach{Matt.22.9}{Matt}{true}\chronologyannotationclaim{composition.gospel-of-matthew}{composition}{catholic-traditional-v1}{disputed}{about A.D. 38-45}{c.~A.D.~38--45}; \chronologyannotationreach{Matt.22.1}{Matt}{true}\chronologyannotationreach{Matt.22.10}{Matt}{true}\chronologyannotationreach{Matt.22.11}{Matt}{true}\chronologyannotationreach{Matt.22.12}{Matt}{true}\chronologyannotationreach{Matt.22.13}{Matt}{true}\chronologyannotationreach{Matt.22.14}{Matt}{true}\chronologyannotationreach{Matt.22.2}{Matt}{true}\chronologyannotationreach{Matt.22.3}{Matt}{true}\chronologyannotationreach{Matt.22.4}{Matt}{true}\chronologyannotationreach{Matt.22.5}{Matt}{true}\chronologyannotationreach{Matt.22.6}{Matt}{true}\chronologyannotationreach{Matt.22.7}{Matt}{true}\chronologyannotationreach{Matt.22.8}{Matt}{true}\chronologyannotationreach{Matt.22.9}{Matt}{true}\chronologyannotationclaim{composition.gospel-of-matthew}{composition}{catholic-traditional-v1}{disputed}{about the year 40-42}{c.~A.D.~40--42}; \chronologyannotationreach{Matt.22.1}{Matt}{true}\chronologyannotationreach{Matt.22.10}{Matt}{true}\chronologyannotationreach{Matt.22.11}{Matt}{true}\chronologyannotationreach{Matt.22.12}{Matt}{true}\chronologyannotationreach{Matt.22.13}{Matt}{true}\chronologyannotationreach{Matt.22.14}{Matt}{true}\chronologyannotationreach{Matt.22.2}{Matt}{true}\chronologyannotationreach{Matt.22.3}{Matt}{true}\chronologyannotationreach{Matt.22.4}{Matt}{true}\chronologyannotationreach{Matt.22.5}{Matt}{true}\chronologyannotationreach{Matt.22.6}{Matt}{true}\chronologyannotationreach{Matt.22.7}{Matt}{true}\chronologyannotationreach{Matt.22.8}{Matt}{true}\chronologyannotationreach{Matt.22.9}{Matt}{true}\chronologyannotationclaim{composition.gospel-of-matthew}{composition}{catholic-traditional-v1}{disputed}{the years 40-45}{A.D.~40--45}; \chronologyannotationreach{Matt.22.1}{Matt}{true}\chronologyannotationreach{Matt.22.10}{Matt}{true}\chronologyannotationreach{Matt.22.11}{Matt}{true}\chronologyannotationreach{Matt.22.12}{Matt}{true}\chronologyannotationreach{Matt.22.13}{Matt}{true}\chronologyannotationreach{Matt.22.14}{Matt}{true}\chronologyannotationreach{Matt.22.2}{Matt}{true}\chronologyannotationreach{Matt.22.3}{Matt}{true}\chronologyannotationreach{Matt.22.4}{Matt}{true}\chronologyannotationreach{Matt.22.5}{Matt}{true}\chronologyannotationreach{Matt.22.6}{Matt}{true}\chronologyannotationreach{Matt.22.7}{Matt}{true}\chronologyannotationreach{Matt.22.8}{Matt}{true}\chronologyannotationreach{Matt.22.9}{Matt}{true}\chronologyannotationclaim{composition.gospel-of-matthew}{composition}{catholic-traditional-v1}{disputed}{about the year 60-68}{c.~A.D.~60--68}; \chronologyannotationreach{Matt.22.1}{Matt}{true}\chronologyannotationreach{Matt.22.10}{Matt}{true}\chronologyannotationreach{Matt.22.11}{Matt}{true}\chronologyannotationreach{Matt.22.12}{Matt}{true}\chronologyannotationreach{Matt.22.13}{Matt}{true}\chronologyannotationreach{Matt.22.14}{Matt}{true}\chronologyannotationreach{Matt.22.2}{Matt}{true}\chronologyannotationreach{Matt.22.3}{Matt}{true}\chronologyannotationreach{Matt.22.4}{Matt}{true}\chronologyannotationreach{Matt.22.5}{Matt}{true}\chronologyannotationreach{Matt.22.6}{Matt}{true}\chronologyannotationreach{Matt.22.7}{Matt}{true}\chronologyannotationreach{Matt.22.8}{Matt}{true}\chronologyannotationreach{Matt.22.9}{Matt}{true}\chronologyannotationclaim{composition.gospel-of-matthew}{composition}{catholic-traditional-v1}{disputed}{about the years 64-67}{c.~A.D.~64--67}; \chronologyannotationreach{Matt.22.1}{Matt}{true}\chronologyannotationreach{Matt.22.10}{Matt}{true}\chronologyannotationreach{Matt.22.11}{Matt}{true}\chronologyannotationreach{Matt.22.12}{Matt}{true}\chronologyannotationreach{Matt.22.13}{Matt}{true}\chronologyannotationreach{Matt.22.14}{Matt}{true}\chronologyannotationreach{Matt.22.2}{Matt}{true}\chronologyannotationreach{Matt.22.3}{Matt}{true}\chronologyannotationreach{Matt.22.4}{Matt}{true}\chronologyannotationreach{Matt.22.5}{Matt}{true}\chronologyannotationreach{Matt.22.6}{Matt}{true}\chronologyannotationreach{Matt.22.7}{Matt}{true}\chronologyannotationreach{Matt.22.8}{Matt}{true}\chronologyannotationreach{Matt.22.9}{Matt}{true}\chronologyannotationclaim{composition.gospel-of-matthew}{composition}{catholic-traditional-v1}{disputed}{about the year 50}{c.~A.D.~50}\\[.3em]%
}
\expandafter\def\csname triptychchronologyorientation@row-31b560fba84818dc\endcsname#1{%
Gospel (shorter form): Matt. 22:1--10 & Narrated event: The parable of the marriage feast, proposed in the Temple & \chronologyannotationreach{Matt.22.1}{Matt.22.1-14}{false}\chronologyannotationreach{Matt.22.10}{Matt.22.1-14}{false}\chronologyannotationreach{Matt.22.2}{Matt.22.1-14}{false}\chronologyannotationreach{Matt.22.3}{Matt.22.1-14}{false}\chronologyannotationreach{Matt.22.4}{Matt.22.1-14}{false}\chronologyannotationreach{Matt.22.5}{Matt.22.1-14}{false}\chronologyannotationreach{Matt.22.6}{Matt.22.1-14}{false}\chronologyannotationreach{Matt.22.7}{Matt.22.1-14}{false}\chronologyannotationreach{Matt.22.8}{Matt.22.1-14}{false}\chronologyannotationreach{Matt.22.9}{Matt.22.1-14}{false}\chronologyannotationclaim{life-of-christ.parable-of-the-marriage-feast}{narrated-event}{catholic-traditional-v1}{preferred}{derived A.D. 29 from February, A.U.C. 782- Passover, 782}{A.D.~29 (derived)}\\*[.1em]%
}
\expandafter\def\csname triptychchronologyorientation@row-84e5dfa573250edb\endcsname#1{%
 & Composition: Matthew; Semitic original associated with Palestine & Disputed: \chronologyannotationreach{Matt.22.1}{Matt}{true}\chronologyannotationreach{Matt.22.10}{Matt}{true}\chronologyannotationreach{Matt.22.2}{Matt}{true}\chronologyannotationreach{Matt.22.3}{Matt}{true}\chronologyannotationreach{Matt.22.4}{Matt}{true}\chronologyannotationreach{Matt.22.5}{Matt}{true}\chronologyannotationreach{Matt.22.6}{Matt}{true}\chronologyannotationreach{Matt.22.7}{Matt}{true}\chronologyannotationreach{Matt.22.8}{Matt}{true}\chronologyannotationreach{Matt.22.9}{Matt}{true}\chronologyannotationclaim{composition.gospel-of-matthew}{composition}{catholic-traditional-v1}{disputed}{about A.D. 38-45}{c.~A.D.~38--45}; \chronologyannotationreach{Matt.22.1}{Matt}{true}\chronologyannotationreach{Matt.22.10}{Matt}{true}\chronologyannotationreach{Matt.22.2}{Matt}{true}\chronologyannotationreach{Matt.22.3}{Matt}{true}\chronologyannotationreach{Matt.22.4}{Matt}{true}\chronologyannotationreach{Matt.22.5}{Matt}{true}\chronologyannotationreach{Matt.22.6}{Matt}{true}\chronologyannotationreach{Matt.22.7}{Matt}{true}\chronologyannotationreach{Matt.22.8}{Matt}{true}\chronologyannotationreach{Matt.22.9}{Matt}{true}\chronologyannotationclaim{composition.gospel-of-matthew}{composition}{catholic-traditional-v1}{disputed}{about the year 40-42}{c.~A.D.~40--42}; \chronologyannotationreach{Matt.22.1}{Matt}{true}\chronologyannotationreach{Matt.22.10}{Matt}{true}\chronologyannotationreach{Matt.22.2}{Matt}{true}\chronologyannotationreach{Matt.22.3}{Matt}{true}\chronologyannotationreach{Matt.22.4}{Matt}{true}\chronologyannotationreach{Matt.22.5}{Matt}{true}\chronologyannotationreach{Matt.22.6}{Matt}{true}\chronologyannotationreach{Matt.22.7}{Matt}{true}\chronologyannotationreach{Matt.22.8}{Matt}{true}\chronologyannotationreach{Matt.22.9}{Matt}{true}\chronologyannotationclaim{composition.gospel-of-matthew}{composition}{catholic-traditional-v1}{disputed}{the years 40-45}{A.D.~40--45}; \chronologyannotationreach{Matt.22.1}{Matt}{true}\chronologyannotationreach{Matt.22.10}{Matt}{true}\chronologyannotationreach{Matt.22.2}{Matt}{true}\chronologyannotationreach{Matt.22.3}{Matt}{true}\chronologyannotationreach{Matt.22.4}{Matt}{true}\chronologyannotationreach{Matt.22.5}{Matt}{true}\chronologyannotationreach{Matt.22.6}{Matt}{true}\chronologyannotationreach{Matt.22.7}{Matt}{true}\chronologyannotationreach{Matt.22.8}{Matt}{true}\chronologyannotationreach{Matt.22.9}{Matt}{true}\chronologyannotationclaim{composition.gospel-of-matthew}{composition}{catholic-traditional-v1}{disputed}{about the year 60-68}{c.~A.D.~60--68}; \chronologyannotationreach{Matt.22.1}{Matt}{true}\chronologyannotationreach{Matt.22.10}{Matt}{true}\chronologyannotationreach{Matt.22.2}{Matt}{true}\chronologyannotationreach{Matt.22.3}{Matt}{true}\chronologyannotationreach{Matt.22.4}{Matt}{true}\chronologyannotationreach{Matt.22.5}{Matt}{true}\chronologyannotationreach{Matt.22.6}{Matt}{true}\chronologyannotationreach{Matt.22.7}{Matt}{true}\chronologyannotationreach{Matt.22.8}{Matt}{true}\chronologyannotationreach{Matt.22.9}{Matt}{true}\chronologyannotationclaim{composition.gospel-of-matthew}{composition}{catholic-traditional-v1}{disputed}{about the years 64-67}{c.~A.D.~64--67}; \chronologyannotationreach{Matt.22.1}{Matt}{true}\chronologyannotationreach{Matt.22.10}{Matt}{true}\chronologyannotationreach{Matt.22.2}{Matt}{true}\chronologyannotationreach{Matt.22.3}{Matt}{true}\chronologyannotationreach{Matt.22.4}{Matt}{true}\chronologyannotationreach{Matt.22.5}{Matt}{true}\chronologyannotationreach{Matt.22.6}{Matt}{true}\chronologyannotationreach{Matt.22.7}{Matt}{true}\chronologyannotationreach{Matt.22.8}{Matt}{true}\chronologyannotationreach{Matt.22.9}{Matt}{true}\chronologyannotationclaim{composition.gospel-of-matthew}{composition}{catholic-traditional-v1}{disputed}{about the year 50}{c.~A.D.~50}\\[.3em]%
}
\expandafter\def\csname triptychchronologyorientation@row-55584471684ef977\endcsname#1{%
Alleluia Verse: Cf. Ephesians 1:17--18 & Composition: Ephesians; Rome or Caesarea; recipients in Asia Minor & Disputed: \chronologyannotationreach{Eph.1.17}{Eph}{true}\chronologyannotationreach{Eph.1.18}{Eph}{true}\chronologyannotationclaim{composition.epistle-to-the-ephesians}{composition}{catholic-traditional-v1}{disputed}{a period between 58 and 63}{A.D.~58--63}; \chronologyannotationreach{Eph.1.17}{Eph}{true}\chronologyannotationreach{Eph.1.18}{Eph}{true}\chronologyannotationclaim{composition.epistle-to-the-ephesians}{composition}{catholic-traditional-v1}{disputed}{(Philemon; Colossians; Ephesians; Philippians), 61}{A.D.~61}\\[.3em]%
}
\expandafter\def\csname triptychchronologyorientation@row-9c28d37a305c76e1\endcsname#1{%
Second Reading: Phil. 4:12--14, 19--20 & Composition: Philippians; Rome to Philippi in the cited account & Disputed: \chronologyannotationreach{Phil.4.12}{Phil}{true}\chronologyannotationreach{Phil.4.13}{Phil}{true}\chronologyannotationreach{Phil.4.14}{Phil}{true}\chronologyannotationreach{Phil.4.19}{Phil}{true}\chronologyannotationreach{Phil.4.20}{Phil}{true}\chronologyannotationclaim{composition.epistle-to-the-philippians}{composition}{catholic-traditional-v1}{disputed}{(Philemon; Colossians; Ephesians; Philippians), 61}{A.D.~61}; \chronologyannotationreach{Phil.4.12}{Phil}{true}\chronologyannotationreach{Phil.4.13}{Phil}{true}\chronologyannotationreach{Phil.4.14}{Phil}{true}\chronologyannotationreach{Phil.4.19}{Phil}{true}\chronologyannotationreach{Phil.4.20}{Phil}{true}\chronologyannotationclaim{composition.epistle-to-the-philippians}{composition}{catholic-traditional-v1}{disputed}{at Rome (A.D. 62-64)}{A.D.~62--64}\\[.3em]%
}
\expandafter\def\csname triptychchronologyorientation@row-c740b284c4f8a93f\endcsname#1{%
Communion Antiphon (alternative): 1 John. 3:2 & Composition: First John; Ephesus probable & \chronologyannotationreach{1John.3.2}{1John}{true}\chronologyannotationclaim{composition.first-epistle-of-st-john}{composition}{catholic-traditional-v1}{preferred}{from the year 90 to 100 (approximately)}{c.~A.D.~90--100}\\[.3em]%
}

% Generated. Do not edit.
% Complete orientation-v3 table; wording, dates and order belong to the corpus.
\begin{chronologysharedtable}
\chronologyorientationrow{row-216603097019c7fc}{}
\chronologyorientationrow{row-3bf36d9c23635597}{}
\chronologyorientationrow{row-41db3bd24cba735a}{}
\midrule
\chronologyorientationrow{row-7de60a4121289421}{}
\chronologyorientationrow{row-1cfd5f3c06636c02}{}
\chronologyorientationrow{row-a9cb00f8a328fe9a}{}
\midrule
\chronologyorientationrow{row-0c09cf47b3a0c749}{}
\chronologyorientationrow{row-410350b07d43c785}{}
\midrule
\chronologyorientationrow{row-b81bab3a554aee2e}{}
\midrule
\chronologyorientationrow{row-bbdb0e76aa03a178}{}
\chronologyorientationrow{row-4d9bb5eddc53b491}{}
\midrule
\chronologyorientationrow{row-31b560fba84818dc}{}
\chronologyorientationrow{row-84e5dfa573250edb}{}
\midrule
\chronologyorientationrow{row-55584471684ef977}{}
\midrule
\chronologyorientationrow{row-9c28d37a305c76e1}{}
\midrule
\chronologyorientationrow{row-c740b284c4f8a93f}{}
\end{chronologysharedtable}
{\footnotesize Psalm numbering: Vg. = Vulgate; Heb. = Hebrew.\par}


\subsection{Evidence for each date}
\label{sec:date-location-evidence}
{\small For each passage, as the Scriptural Date and Location table names it: what is dated, how the corpus weighs each answer, the source's own words, and the sources cited. A subject the corpus does not date is reported as such and given no date.\par}
{\small
\par\smallskip\noindent\textit{Responsorial Psalm: Ps. 23:1--3a, 3b--4, 5, 6 (Heb.; Vg. 22), Ps. 23:6cd (Heb.; Vg. 22)}\par\nobreak
\begin{itemize}[leftmargin=1.2em,itemsep=1pt,topsep=1pt,parsep=0pt]
\item \textbf{Historical setting}. No date: the chronology corpus has not yet dated this (research pending).
\item \textbf{Composition} (preferred). Psalms: Before c.~165 B.C. (``before the Maccabean period, around 165 B.C.; no individual psalm can be dated securely''; Catholic critical chronology; New American Bible, Revised Edition --- USCCB web state (2026), Psalms introduction).
\end{itemize}
\par\smallskip\noindent\textit{Communion Antiphon (first): Cf. Psalm 33:11 (Vg.; Heb. 34)}\par\nobreak
\begin{itemize}[leftmargin=1.2em,itemsep=1pt,topsep=1pt,parsep=0pt]
\item \textbf{Traditional attribution} (disputed). David's reign: B.C.~1055--1015 (``reigned from 1055 to 1015 B.C.''; The Catholic Encyclopedia, Volume 4 (New York, 1908), article 'David, King,' printed p. 642, opening paragraph). Place: Hebron, then Jerusalem (The Holy Bible, Douay-Rheims (Challoner revision), Project Gutenberg eBook 1581: 2 Kings 5:5).
\item \textbf{Superscription setting} (preferred). David's flight to Achis, king of Geth: A.M.~2944 (Haydock edition reporting Ussher) (``A. M. 2944, A. C. 1060''; Douay Rheims Edition of the Holy Catholic Bible with a Comprehensive Catholic Commentary, William Peter Feeney 2014 Memorial Edition, 1K.21.1-15; Douay Rheims Edition of the Holy Catholic Bible with a Comprehensive Catholic Commentary, William Peter Feeney 2014 Memorial Edition, Ps.70.1). Place: Geth, at the Philistine court (The Holy Bible, Douay-Rheims (Challoner revision), Project Gutenberg eBook 1581: Ps. 33:1, 1 Kings 21:10, 21:15).
\item \textbf{Composition}. As under Responsorial Psalm: Ps. 23:1--3a, 3b--4, 5, 6 (Heb.; Vg. 22), Ps. 23:6cd (Heb.; Vg. 22).
\end{itemize}
\par\smallskip\noindent\textit{Entrance Antiphon: Ps. 129:3--4 (Vg.; Heb. 130)}\par\nobreak
\begin{itemize}[leftmargin=1.2em,itemsep=1pt,topsep=1pt,parsep=0pt]
\item \textbf{Historical setting}. As under Responsorial Psalm: Ps. 23:1--3a, 3b--4, 5, 6 (Heb.; Vg. 22), Ps. 23:6cd (Heb.; Vg. 22).
\item \textbf{Composition}. As under Responsorial Psalm: Ps. 23:1--3a, 3b--4, 5, 6 (Heb.; Vg. 22), Ps. 23:6cd (Heb.; Vg. 22).
\end{itemize}
\par\smallskip\noindent\textit{First Reading: Is. 25:6--10a}\par\nobreak
\begin{itemize}[leftmargin=1.2em,itemsep=1pt,topsep=1pt,parsep=0pt]
\item \textbf{Traditional attribution} (preferred). Isaias (ministry in Souvay's traditional account): B.C.~740--701 (``from the closing year of Ozias, King of Juda (740 B.C.); no prophecies appear to be later than 701 B.C.''; The Catholic Encyclopedia, Volume 8 (New York, 1910)). Place: Judah and Jerusalem (The Holy Bible, Douay-Rheims (Challoner revision), Project Gutenberg eBook 1581: Is. 1:1).
\end{itemize}
\par\smallskip\noindent\textit{Gospel (longer form): Matt. 22:1--14}\par\nobreak
\begin{itemize}[leftmargin=1.2em,itemsep=1pt,topsep=1pt,parsep=0pt]
\item \textbf{Narrated event} (preferred). The parable of the marriage feast, proposed in the Temple: A.D.~29 (derived) (``derived A.D. 29 from February, A.U.C. 782- Passover, 782''; The Catholic Encyclopedia, Volume 8 (New York, 1910), article 'Jesus Christ,' III.B.3, printed p. 378 and pp. 380-381: chronological calibration and Sixth, Seventh, and Ninth Journeys).
\item \textbf{Composition} (disputed). The Gospel of St Matthew: c.~A.D.~38--45 (``about A.D. 38-45''; The Catholic Encyclopedia, Volume 10 (1911); The Catholic Encyclopedia, Volume 14 (New York, 1912)); c.~A.D.~40--42 (``about the year 40-42''; The Catholic Encyclopedia, Volume 10 (1911)); A.D.~40--45 (``the years 40-45''; The Catholic Encyclopedia, Volume 10 (1911)); c.~A.D.~60--68 (``about the year 60-68''; The Catholic Encyclopedia, Volume 10 (1911)); c.~A.D.~64--67 (``about the years 64-67''; The Catholic Encyclopedia, Volume 10 (1911)); c.~A.D.~50 (``about the year 50''; The Catholic Encyclopedia, Volume 14 (New York, 1912)). Place: Semitic original associated with Palestine (The Catholic Encyclopedia, Volume 10 (1911)).
\end{itemize}
\par\smallskip\noindent\textit{Gospel (shorter form): Matt. 22:1--10}\par\nobreak
\begin{itemize}[leftmargin=1.2em,itemsep=1pt,topsep=1pt,parsep=0pt]
\item \textbf{Narrated event}. As under Gospel (longer form): Matt. 22:1--14.
\item \textbf{Composition}. As under Gospel (longer form): Matt. 22:1--14.
\end{itemize}
\par\smallskip\noindent\textit{Alleluia Verse: Cf. Ephesians 1:17--18}\par\nobreak
\begin{itemize}[leftmargin=1.2em,itemsep=1pt,topsep=1pt,parsep=0pt]
\item \textbf{Composition} (disputed). The Epistle to the Ephesians: A.D.~58--63 (``a period between 58 and 63''; The Catholic Encyclopedia, Volume 5 (New York, 1909)); A.D.~61 (``(Philemon; Colossians; Ephesians; Philippians), 61''; The Catholic Encyclopedia, Volume 11 (New York, 1911)). Place: Rome or Caesarea; recipients in Asia Minor (The Catholic Encyclopedia, Volume 5 (New York, 1909)).
\end{itemize}
\par\smallskip\noindent\textit{Second Reading: Phil. 4:12--14, 19--20}\par\nobreak
\begin{itemize}[leftmargin=1.2em,itemsep=1pt,topsep=1pt,parsep=0pt]
\item \textbf{Composition} (disputed). The Epistle to the Philippians: A.D.~61 (``(Philemon; Colossians; Ephesians; Philippians), 61''; The Catholic Encyclopedia, Volume 11 (New York, 1911)); A.D.~62--64 (``at Rome (A.D. 62-64)''; The Catholic Encyclopedia, Volume 12 (1911)). Place: Rome to Philippi in the cited account (The Catholic Encyclopedia, Volume 12 (1911); The Holy Bible, Douay-Rheims (Challoner revision), Project Gutenberg eBook 1581: Phil. 1:1).
\end{itemize}
\par\smallskip\noindent\textit{Communion Antiphon (alternative): 1 John. 3:2}\par\nobreak
\begin{itemize}[leftmargin=1.2em,itemsep=1pt,topsep=1pt,parsep=0pt]
\item \textbf{Composition} (preferred). The First Epistle of St John: c.~A.D.~90--100 (``from the year 90 to 100 (approximately)''; The Catholic Encyclopedia, Volume 14 (New York, 1912); The Catholic Encyclopedia, Volume 8 (New York, 1910)). Place: Ephesus probable (The Catholic Encyclopedia, Volume 8 (New York, 1910)).
\end{itemize}
\par}
